\documentclass[10pt,a4paper]{amsart}
\usepackage[margin=19mm]{geometry}
\usepackage{amsmath,amssymb,mathtools}
\usepackage[scr=rsfs,cal=euler]{mathalfa}
\usepackage{xfrac}
\usepackage[unicode,hidelinks,bookmarks=false]{hyperref}
\newcommand{\ie}{i.e.\ }
\newcommand{\cf}{cf.\ }
\newcommand{\resp}{respectively\ }
\newcommand{\Subsection}{\subsection}
\providecommand{\nobreakdash}{\nobreak\hbox{-}}
\setlength{\emergencystretch}{2em}
\multlinegap0pt
\usepackage{mathtools}
\usepackage{amssymb}
\usepackage[shortlabels]{enumitem}
\usepackage{tikz}
\usepackage{float}
\newtheorem{theorem}{Theorem}
\newcommand{\N}{\mathbb{N}}
\title{\textbf{Recurrence Structures, Finite State Decomposition, and Statistical Bias in Collatz Path Sequences}}
\author{Sawon Pratiher}
\date{Source: arXiv:1608.03600v7 (CC BY 4.0)}
\begin{document}
\maketitle

\noindent\emph{Source and licence.} \textbf{Recurrence Structures, Finite State Decomposition, and Statistical Bias in Collatz Path Sequences}. Version arXiv:1608.03600v7 (CC BY 4.0), distributed by arXiv under the Creative Commons Attribution 4.0 International licence, http://creativecommons.org/licenses/by/4.0/, as recorded in the arXiv licence metadata for this exact version and shown on the abstract page https://arxiv.org/abs/1608.03600v7. Full licence notice: \texttt{licenses/arxiv-1608.03600-CC-BY-4.0.txt}.\par\smallskip

\noindent\textit{Editorial scope.} Counted unit: the article's theorem thm:fsm\_transitions (source0.tex lines 239--284). The article's complete original proof of it is delivered with it (source0.tex lines 286--294, one \texttt{proof} environment of 9 lines). No mathematics is written, repaired or re-derived: the statement and the proof are the article's own lines, in the article's own notation, and no retained line is substituted or re-flowed.  No further source-local context is needed: every cross-reference of every retained line resolves inside the two delivered spans.  Nothing was omitted from any retained span, so the package carries no omission record.  No retained line contains a citation, so the wrapper carries no bibliography and no bibliography entry.  No retained line contains a figure environment, an image-inclusion command or drawing code, so no image is delivered and the package declares no asset dependency.  Editorial formatting only, in addition: an AMS \texttt{amsart} preamble that restates the article's own declarations and macros used by the retained lines, namely the environments \texttt{theorem} on separate counters, together with the standard packages those lines need.  The retained environments print in this document as Theorem 1 in order of appearance, whereas the article prints the same items with the numbering of its own full document.  The complete native source of the article as distributed in the arXiv e-print for this exact version is bundled unchanged as \texttt{sources/100626/source0.tex}.
\begin{theorem}[FSM transitions]\label{thm:fsm_transitions}
  Define the transition function $\delta: \{\mathfrak{a}, \mathfrak{b}, \mathfrak{c}, \mathfrak{d}, \mathfrak{e}, \mathfrak{f}\} \times \N_0 \to \{\mathfrak{a}, \mathfrak{b}, \mathfrak{c}, \mathfrak{d}, \mathfrak{e}, \mathfrak{f}\} \times \N_0$ by the following rules, where each transition is realized by a specific finite sequence of Collatz operations, and $n$ denotes the current parameter:

  \medskip
  \noindent\textbf{(i) From form $\mathfrak{a}$: $9n+8$.}
  \begin{align}
    n \text{ even} &: \quad 9n+8 \xrightarrow{E} 9\!\left(\tfrac{n}{2}\right)+4, && \text{transition } \mathfrak{a} \to \mathfrak{b}, \text{ parameter } n \mapsto \tfrac{n}{2}. \label{eq:a_even} \\
    n \text{ odd} &: \quad 9n+8 \xrightarrow{O,E} 9\!\left(\tfrac{3n+1}{2}\right)+8, && \text{transition } \mathfrak{a} \to \mathfrak{a}, \text{ parameter } n \mapsto \tfrac{3n+1}{2}. \label{eq:a_odd}
  \end{align}

  \medskip
  \noindent\textbf{(ii) From form $\mathfrak{b}$: $9n+4$.}
  \begin{align}
    n \text{ even} &: \quad 9n+4 \xrightarrow{E} 9\!\left(\tfrac{n}{4}\right)+2, && \text{transition } \mathfrak{b} \to \mathfrak{c}, \text{ parameter } n \mapsto \tfrac{n}{4}. \label{eq:b_even} \\
    n \text{ odd} &: \quad 9n+4 \xrightarrow{O,E} 9\!\left(\tfrac{3n+1}{4}\right)+2, && \text{transition } \mathfrak{b} \to \mathfrak{c}, \text{ parameter } n \mapsto \tfrac{3n+1}{4}. \label{eq:b_odd}
  \end{align}
  (Note: in the even case, the parameter $n$ must satisfy $4 \mid n$; the intermediate cases when $n \equiv 2 \pmod{4}$ are handled by further parity analysis.)

  \medskip
  \noindent\textbf{(iii) From form $\mathfrak{c}$: $9n+2$.}
  \begin{align}
    n \text{ even} &: \quad 9n+2 \xrightarrow{E} 9\!\left(\tfrac{n}{8}\right)+1, && \text{transition } \mathfrak{c} \to \mathfrak{d}, \text{ parameter } n \mapsto \tfrac{n}{8}. \label{eq:c_even} \\
    n \text{ odd} &: \quad 9n+2 \xrightarrow{O,E} 9\!\left(\tfrac{3n-1}{2}\right)+8, && \text{transition } \mathfrak{c} \to \mathfrak{a}, \text{ parameter } n \mapsto \tfrac{3n-1}{2}. \label{eq:c_odd}
  \end{align}

  \medskip
  \noindent\textbf{(iv) From form $\mathfrak{d}$: $9n+1$.}
  \begin{align}
    n \text{ even} &: \quad 9n+1 \xrightarrow{E} 9\!\left(\tfrac{3n}{2}\right)+2, && \text{transition } \mathfrak{d} \to \mathfrak{c}, \text{ parameter via intermediate steps.} \label{eq:d_even} \\
    n \text{ odd} &: \quad 9n+1 \xrightarrow{E} 9\!\left(\tfrac{n-1}{2}\right)+5, && \text{transition } \mathfrak{d} \to \mathfrak{e}, \text{ parameter } n \mapsto \tfrac{n-1}{2}. \label{eq:d_odd}
  \end{align}

  \medskip
  \noindent\textbf{(v) From form $\mathfrak{e}$: $9n+5$.}
  \begin{align}
    n \text{ even} &: \quad 9n+5 \xrightarrow{O,E} 9\!\left(\tfrac{3n}{2}\right)+8, && \text{transition } \mathfrak{e} \to \mathfrak{a}, \text{ parameter } n \mapsto \tfrac{3n}{2}. \label{eq:e_even} \\
    n \text{ odd} &: \quad 9n+5 \xrightarrow{E} 9\!\left(\tfrac{n-1}{2}\right)+7, && \text{transition } \mathfrak{e} \to \mathfrak{f}, \text{ parameter } n \mapsto \tfrac{n-1}{2}. \label{eq:e_odd}
  \end{align}

  \medskip
  \noindent\textbf{(vi) From form $\mathfrak{f}$: $9n+7$.}
  \begin{align}
    n \text{ even} &: \quad 9n+7 \xrightarrow{O,E,E} 9\!\left(\tfrac{3n+1}{4}\right)+8, && \text{transition } \mathfrak{f} \to \mathfrak{a}, \text{ parameter } n \mapsto \tfrac{3n+1}{4}. \label{eq:f_even} \\
    n \text{ odd} &: \quad 9n+7 \xrightarrow{O,E,E} 9\!\left(\tfrac{3n+1}{4}\right)+2, && \text{transition } \mathfrak{f} \to \mathfrak{c}, \text{ parameter via further analysis.} \label{eq:f_odd}
  \end{align}
\end{theorem}

\begin{proof}
  We verify each case by direct algebraic computation. We illustrate with the transitions from form~$\mathfrak{a}$.

  \textbf{Case (i), $n$ even:} Let $n = 2m$. Then $9n+8 = 18m+8$, which is even. Applying $E$: $(18m+8)/2 = 9m+4 = 9(\tfrac{n}{2})+4$, which has form~$\mathfrak{b}$ with parameter $\tfrac{n}{2} = m$.

  \textbf{Case (i), $n$ odd:} Let $n = 2m+1$. Then $9n+8 = 18m+17$, which is odd. Applying $O$: $3(18m+17)+1 = 54m+52$, which is even. Applying $E$: $(54m+52)/2 = 27m+26 = 9(3m+2)+8$. Since $n = 2m+1$, we have $\tfrac{3n+1}{2} = \tfrac{3(2m+1)+1}{2} = 3m+2$, confirming form~$\mathfrak{a}$ with parameter $\tfrac{3n+1}{2}$.

  The remaining cases follow by analogous computations, tracking the Collatz operations and verifying the resulting residue modulo~9.
\end{proof}

\end{document}
